\section{Mathematics}\label{sec:math}

A common way to move \LaTeX{} mathematics into Word is to paste each equation in as a
picture or to flatten it to plain text. \texdocx{} does neither. Its math parser reads
the same expanded token stream as the rest of the document, so your own macros work inside
formulas, and it recovers the \emph{structure} of each formula (fractions, scripts, radicals,
delimiters, matrices) rather than its layout. The \texttt{docx} writer then turns that
structure into Office Math Markup Language (OMML), the format Word's own equation editor
uses. Every formula in this section is therefore a native Word equation: click on it and
you can edit it, change its font size, or copy it into another document.

\subsection{Inline and displayed formulas}

Inline math such as $e^{i\pi} + 1 = 0$ or $\R^n$ stays in the line of text, just as it does in
\LaTeX. Fractions like $\frac{a}{b}$, roots like $\sqrt{2}$ and Greek letters such as
$\alpha$, $\beta$ and $\Gamma$ all map to their OMML counterparts. A displayed equation
gets its own paragraph. The Gaussian integral is a good first example:
\begin{equation}\label{eq:math-gauss}
  \int_{-\infty}^{\infty} e^{-x^2} \, dx = \sqrt{\pi}.
\end{equation}
Because the integral has its limits attached as scripts, Word places them beside the sign,
which is also what \LaTeX{} does for \cmd{int}. Sums and products put their limits above
and below in display style:
\[
  \sum_{k=1}^{n} k = \frac{n(n+1)}{2},
  \qquad
  \lim_{n \to \infty} \left(1 + \frac{1}{n}\right)^{n} = e.
\]
Here \cmd{left} and \cmd{right} become a single OMML delimiter object, so Word grows the
parentheses with their contents. The limit under \cmd{lim} is a real lower-limit object, not a
subscript.

\subsection{Numbering and cross-references}

Equation~\eqref{eq:math-gauss} above was numbered automatically. In the Word file the number
is not plain text: it is a \texttt{SEQ Equation} field, wrapped in a bookmark named after the
\cmd{label}. The equation itself sits at a centre tab stop and the number at a right tab stop
on the same line. Each \cmd{eqref} becomes a \texttt{REF} field that points at that bookmark,
with the parentheses added around it, exactly like the reference you just read. Both fields
carry the number \texdocx{} computed as their cached value, so the document is correct as soon
as it opens; the fields are designed so that updating them in Word (F9) gives the same numbers.

Multi-line environments number each row separately. In an \env{align} environment every row
becomes its own Word paragraph, so that it can carry its own number, and \cmd{notag} (or
\cmd{nonumber}) suppresses the number of a single row:
\begin{align}
  f(x) &= (x + 1)^2 \label{eq:math-square} \\
       &= x^2 + 2x + 1 \notag \\
  f'(x) &= 2x + 2. \label{eq:math-derivative}
\end{align}
The middle row has no number, so the derivative in~\eqref{eq:math-derivative} directly follows
\eqref{eq:math-square}, just as it would in a PDF. One visible difference remains: because the rows are separate paragraphs,
Word does not yet line up the equals signs across them. In the starred forms, which
need no numbers, the rows stay in one OMML equation array and the \& marks do align.
The starred forms number nothing; here is
\env{gather*}, which centres each row on its own:
\begin{gather*}
  \binom{n}{k} = \frac{n!}{k!\,(n-k)!}, \\
  \sqrt[3]{27} = 3, \qquad \sqrt{x^2 + y^2} \geq 0.
\end{gather*}
Note that \cmd{binom} becomes a bar-less fraction inside parentheses, and the optional argument
of \cmd{sqrt} becomes the degree of an OMML radical.

\subsection{Matrices and cases}

The \env{pmatrix} and \env{bmatrix} environments become OMML matrices wrapped in round or
square delimiters:
\[
  A = \begin{pmatrix} 1 & 2 \\ 3 & 4 \end{pmatrix},
  \qquad
  I_3 = \begin{bmatrix} 1 & 0 & 0 \\ 0 & 1 & 0 \\ 0 & 0 & 1 \end{bmatrix}.
\]
A \env{cases} environment becomes a left-aligned two-column matrix behind a single opening
brace, with no closing delimiter:
\begin{equation}\label{eq:math-abs}
  |x| = \begin{cases}
    x  & \text{if}\ x \geq 0, \\
    -x & \text{otherwise.}
  \end{cases}
\end{equation}
Text inside \cmd{text} is written as normal (non-math) text in the equation, so it keeps
upright letters and its spaces, including those at the start or end of the argument, so
\verb|\text{ if }| reads exactly as it does in the PDF.

\subsection{Fonts, braces and operators}

Math alphabets are mapped to Unicode mathematical letters, which Word shows directly in the
equation:
\cmd{mathbb} gives $\mathbb{N} \subset \mathbb{Z} \subset \mathbb{Q} \subset \R$, and
\cmd{mathcal} gives $\mathcal{F}$, $\mathcal{L}$ and $\mathcal{O}(n \log n)$. The preamble of
this guide defines \cmd{R} and \cmd{E} as shorthands for $\R$ and $\E$, and they behave like
any other macro.

An \cmd{overbrace} with a superscript becomes an OMML group character with an upper limit, so
the label sits above the brace:
\begin{equation}\label{eq:math-brace}
  \overbrace{1 + 1 + \cdots + 1}^{n\ \text{terms}} = n.
\end{equation}
Operators declared with \cmd{DeclareMathOperator} are set upright as a single function name.
This guide's preamble declares \cmd{argmax}, which reads naturally in an expected-value
problem:
\begin{equation}\label{eq:math-argmax}
  \theta^{\ast} = \argmax_{\theta \in \Theta} \E\left[ \log p(x \mid \theta) \right].
\end{equation}
Because the preamble uses the unstarred form, the subscript stays to the side even in display
style, exactly as \LaTeX{} sets it; the starred \cmd{DeclareMathOperator*} would place it
underneath, like \cmd{lim}. Inline, as in $\argmax_{\theta} f(\theta)$, the subscript is
always at the side.

\begin{tip}
The numbers of equations~\eqref{eq:math-gauss} to~\eqref{eq:math-argmax}, and every reference
to them, are fields. If you rearrange equations in Word, select the whole document and press F9
to update the fields: the equation numbers are recounted and each reference follows its
bookmark.
\end{tip}
